Deicing for UAV's most used method:
- we could use the liedat at the fron of the plane to detect ahead of time the icing conditions and either avoid them or initiate the deicing equipemtn

-include an ice detecting sensor mandatory: Fiber-optic Ice Detection,  Heat of Fusion Sensor

-the most common use of deicing methods employ the use of an electric system tha pulses a current through a wire and that heats up the element and melts the ice

-use of coatings might be a better use since we are a little more problematic when it comes tot he desings of the wings

-passive elements will be better for us since we dont need to rebuild the drones but we can coate them with: 
                    -(prevents ice formation)Fluoropolymers applied in various forms,such as ink or aqueous dispersion
                    - (prevents ice stagnation)Liquid-Like Surfaces - These surfaces consist of a thin polymer layer (1 5 nm) characterized by a low CAH (<5°). The low CAH is ascribable to the mobility of the flexible tethered molecules, which is correlated to the length of the polymer chains and the grafting density
                    -(prevents ice stagnation) Phase Change Materials in the Coating System - this one looks the most usefull one since it prevents the droplet from freezing on the surface by absorbing latent hear or giving it out. it can be used in a liquid coating of inert mat
                    - 
-the coating work because we have carbon fibre that can be coated with a lot of materials that are inert and will only serve a transfer liquid transporter -> source this

the rest is the material people's job to do as they have the source (Passive Ice Protection Systems for Unmanned Aerial
Vehicles Applications: A Review)